\documentclass[11pt]{article}
\usepackage{amsmath,amssymb,amsthm}
\usepackage{mathtools}
\usepackage[margin=1.1in]{geometry}
\allowdisplaybreaks
% Theorem environments (shared counter across the whole paper)
\newtheorem{theorem}{Theorem}
\newtheorem{lemma}[theorem]{Lemma}
\newtheorem{corollary}[theorem]{Corollary}
\newtheorem{proposition}[theorem]{Proposition}
\newtheorem{definition}[theorem]{Definition}
\theoremstyle{remark}
\newtheorem{remark}[theorem]{Remark}
\newtheorem{example}[theorem]{Example}
% Macro union across the assembled notes
\newcommand{\vol}{\operatorname{vol}}
\newcommand{\diam}{\operatorname{diam}}
\newcommand{\sgn}{\operatorname{sgn}}
\newcommand{\Hess}{\operatorname{Hess}}
\newcommand{\Var}{\operatorname{Var}}
\newcommand{\dM}{d_M}
\newcommand{\rhoG}{\rho_G}
\newcommand{\Rem}{R}
\newcommand{\NL}{N_\Lambda}
\newcommand{\E}{\mathbb E}
\DeclareMathOperator{\Ric}{Ric}
\DeclareMathOperator{\Tr}{Tr}

\title{Toward the Angular-Preservation Theorem:\\
Reductions, a Non-Radiality Calculus, and the Green-Rank Frontier}
\author{Andrew H. Bond\\
\small Department of Computer Engineering, San Jos\'e State University\\
\small \texttt{andrew.bond@sjsu.edu} \ \textbullet\ \ ORCID 0009-0003-2599-6158}
\date{}

\begin{document}
\maketitle

\begin{abstract}
This is a companion to ``Keep the Angle: A Geometry-Preserving Basis in Spectral
Embeddings'' (Paper~I). Paper~I reduces the rank form of \emph{Angular Preservation}---the
conjecture that row-normalized spectral embeddings keep pairwise angular distances
rank-faithful to geodesic distance---for intrinsic dimension $d\le3$, to two ingredients: a
graph-to-continuum \emph{rank transfer} and \emph{Green-rank positivity}
$\rho_G(M)=\rho_S(d_M(X,Y),-G_M(X,Y))>0$, where $G_M$ is the zero-mean Green kernel. Here we
move both toward theorems. We (i) delete the auxiliary hypotheses (H3) and atomlessness from
Paper~I's conditional results; (ii) prove that the transfer's anti-concentration hypothesis,
false in its bounded-density form, holds in the operative \emph{modulus} form for every
real-analytic metric; (iii) develop a \emph{non-radiality calculus} that reduces Green-rank
positivity to a single scale-invariant index of the Green kernel, via a master inequality
$\rho_G\ge1-12\,\phi\lVert R\rVert_\infty$ and its distributional refinement
$\rho_G\ge1-18(\phi^2V)^{1/3}$; (iv) bound that index geometrically, identifying it with the
\emph{within-shell variance of the heat kernel}---which vanishes on two-point homogeneous
spaces, is finite exactly for $d\le3$, and yields a quantitative $\eta^{2/3}$ perturbative
openness rate---leaving one named curvature-anisotropy amplitude as the sole non-explicit
constant; and (v) carry the whole chain past the critical dimension via a fractional filter,
proving the exact symmetric-space anchor ($\rho_G=1$) in \emph{every} dimension by heat-kernel
subordination. A numerical counterexample hunt across four geometry families (the public
repository) finds none. Green-rank positivity is thereby proved on an increasing hierarchy of
explicit classes and reduced, in general, to one geometric quantity.

\smallskip
\noindent\emph{Verification note.} The results were developed and checked under a
registration-and-independent-verification discipline: every theorem carries a written proof
and a fresh-context adversarial re-derivation, and that process caught and corrected several
errors during preparation. This is not a substitute for human refereeing or machine-checked
proof, and the heat-kernel bound of \S\ref{sec:l2}, in particular, would benefit from an
expert reading; statements are made at the confidence their verification supports.
\end{abstract}

\section{Introduction}\label{sec:intro}

Paper~I (``Keep the Angle'') shows that the geometry a large graph Laplacian appears to lose
to the commute-time collapse survives in the \emph{angular} coordinate of the spectral
embedding, and conjectures \emph{Angular Preservation}: row-normalization keeps pairwise
angular distances rank-faithful to geodesic distance, uniformly in the mode count. For
$d\le3$ the paper reduces the rank form to a chain,
\[
\text{(graph angular ranking)}\ \xRightarrow{\ \text{rank transfer}\ }\
\text{(continuum angular ranking)}\ \xRightarrow{\ \text{Green-kernel rank limit}\ }\
\rho_G(M)>0,
\]
the last step being the Green-kernel rank-limit theorem, which sends the truncated commute
angular ranking to the ranking of the zero-mean Green kernel $G_M$, and leaves open two
things: the graph-level \emph{rank transfer} (governed by a calibration error and an
anti-concentration hypothesis) and \emph{Green-rank positivity},
$\rho_G(M):=\rho_S\bigl(d_M(X,Y),-G_M(X,Y)\bigr)>0$ for $(X,Y)$ i.i.d.\ from the volume.

This companion collects and extends the reductions of both. Its contributions, by part:

\begin{itemize}
\item \textbf{Part~\ref{sec:reductions} (reductions and the transfer hypothesis).} We show
(H3), the injectivity hypothesis of Paper~I's conditional theorem, is not independent---it
follows from (H1)--(H2) under an explicit pinching condition with a quantitative angular gap;
and that the atomlessness assumptions of the rank-limit results are automatic under a bounded
pair density. We then prove the graph-level \emph{rank transfer} needs neither bi-Lipschitz
geometry nor a resolution cutoff---only calibration plus anti-concentration---and rederive
the sharpened growing-window threshold $E_n\sqrt{A_m}\to0$ by an independent route.

\item \textbf{Part~\ref{sec:ac} (the anti-concentration hypothesis).} The transfer's
anti-concentration hypothesis (AC), stated as a bounded density, is \emph{false}---a
nondegenerate saddle forces a logarithmic (van Hove) divergence, guaranteed on any closed
surface of nonpositive Euler characteristic. We isolate the \emph{modulus} form (AC$'$) the
transfer proof actually consumes and prove it for every real-analytic metric, via a uniform
\L ojasiewicz sublevel estimate made uniform over the analytic family by o-minimality.

\item \textbf{Part~\ref{sec:positivity} (a non-radiality calculus).} We reduce Green-rank
positivity to a single quantity. A \emph{master inequality} $\rho_G\ge1-12\,\phi_\Psi
\lVert R\rVert_\infty$ controls the rank deficit by the sup-norm non-radiality of $-G$ against
its best monotone-radial profile; a \emph{distributional refinement} $\rho_G\ge1-18
(\phi_\Psi^2 V)^{1/3}$ replaces the sup by the $L^2$ non-radiality $V$, with the index
$\phi_\Psi^2 V$ dimensionless (scale-invariant)---as a bound on a rank correlation must be.
Two-point homogeneous spaces are the exact ($R\equiv0$) anchor.

\item \textbf{Part~\ref{sec:l2} (the geometric bound).} We identify the $L^2$ non-radiality
with the \emph{within-distance-shell variance of the Green kernel}, and via the heat
representation with the integrated shell-variance of the heat kernel: it vanishes on
two-point homogeneous spaces, is finite \emph{exactly} for $d\le3$ (the same critical
dimension four, now as an integrability threshold), has an explicit spectral-gap tail, and is
$O(\lVert g-g_0\rVert_{C^2}^2)$ near any symmetric metric---so $\rho_G\ge1-C\eta^{2/3}\to1$, a
\emph{quantitative} form of Paper~I's qualitative openness. One named curvature-anisotropy
amplitude remains the sole non-explicit constant.

\item \textbf{Part~\ref{sec:fractional} (past the critical dimension).} For $d\ge4$ the
commute Green kernel leaves $L^2$; a fractional filter $\mu_k^{-s/2}$, $s>d/4$, restores the
whole chain, and heat-kernel subordination proves the exact anchor---fractional Green-rank
$=1$ on all rank-one symmetric spaces---in \emph{every} dimension.
\end{itemize}

The upshot (\S\ref{sec:frontier}): Green-rank positivity is proved exactly on two-point
homogeneous spaces (all dimensions), quantitatively on their $C^2$-neighborhoods, and on the
explicit near-radial / $L^2$-near-radial classes; in general it is reduced to positivity of a
single geometric coefficient, with a numerical counterexample hunt (four families) finding
none. Notation follows Paper~I and its supplement throughout; the phrase ``companion note''
in the assembled parts below refers to the other parts of this paper.

\section{Reductions and the transfer hypothesis}\label{sec:reductions}

\noindent\textbf{Purpose.} This note records four self-contained results that move the
rank form of Angular Preservation (Paper~I, Conjecture on Angular Preservation) toward a
theorem. Notation follows Paper~I and its supplement: $M$ is a closed Riemannian
$d$-manifold with volume normalized to $1$; $F=(\varphi_k/\sqrt{\mu_k})_{k\le m}$ (or its
spectral-cutoff form $F_\Lambda$) is the commute-weighted eigenmap, $R=\lVert F\rVert$,
$N=F/R$; (H1)--(H3), the calibration error $E_n$, the radius band $[\rho^-,\rho^+]$, the
Green-rank coefficient $\rho_G(M)=\rho_S\bigl(d_M(X,Y),-G_M(X,Y)\bigr)$, the diagonal
scale $A_\Lambda$, and the Green-kernel rank-limit theorem are as in Paper~I.
Facts imported as standard are listed in \S5 with sources.

The upshot (\S\ref{sec:frontier}): for $d\le3$, the rank form of Angular Preservation
reduces to
\[
\underbrace{\rho_G(M)>0}_{\text{geometric}}\ +\
\underbrace{\text{(AC)}}_{\text{anti-concentration}}\ +\
\underbrace{E_n\to0\ \ (\text{fixed }m)\ \big/\ E_n\sqrt{A_m}\to0\ \ (\text{growing }m)}_{\text{calibration}},
\]
with Theorem~\ref{thm:open} proving the first item on a $C^2$-open set of metrics,
Theorem~\ref{thm:ac-transfer} proving the transfer given the second and third, and
Lemmas~\ref{lem:h3} and~\ref{lem:atomless} deleting hypotheses ((H3); atomlessness) from
the existing results.

\section{(H3) is not an independent hypothesis}

Recall the setting of Paper~I's conditional theorem: (H1) is $L_0$-bi-Lipschitzness of
$F$ on geodesic $r_0$-balls; (H2) is $\rho^-\le R\le\rho^+$ with $R$ globally
$\Lambda_0$-Lipschitz and $\Lambda_0<1/L_0$; (H3) is injectivity of $N$ on $M$. The key
observation is an exact identity for radial alignment.

\begin{lemma}[Radial-alignment identity]\label{lem:align}
If $N(x)=N(y)$ then $F(x)=tF(y)$ with $t=R(x)/R(y)>0$, and
\[
\lVert F(x)-F(y)\rVert=\lvert t-1\rvert\,R(y)=\lvert R(x)-R(y)\rvert .
\]
\end{lemma}

\begin{proof}
$N(x)=N(y)$ means $F(x)/R(x)=F(y)/R(y)$, i.e.\ $F(x)=tF(y)$ with $t=R(x)/R(y)$. Then
$\lVert F(x)-F(y)\rVert=\lvert t-1\rvert\lVert F(y)\rVert=\lvert tR(y)-R(y)\rvert
=\lvert R(x)-R(y)\rvert$.
\end{proof}

\begin{lemma}[(H3) from (H1)--(H2)]\label{lem:h3}
Assume (H1) on $r_0$-balls and (H2), and set
\[
c_F:=\min\{\lVert F(x)-F(y)\rVert:\ d_M(x,y)\ge r_0\},\qquad
\delta:=\min\bigl(\rho^+-\rho^-,\ \Lambda_0\diam M\bigr).
\]
Then:
\begin{itemize}
\item[(a)] $N$ is injective on every $r_0$-ball (no hypothesis beyond (H1)--(H2));
\item[(b)] if $c_F>\delta$, then (H3) holds on all of $M$, and quantitatively
\[
c_0:=\min\{\lVert N(x)-N(y)\rVert:\ d_M(x,y)\ge r_0\}\ \ge\
\frac{\sqrt{c_F^2-\delta^2}}{\rho^+}\ >\ 0 ;
\]
\item[(c)] if (H1) holds globally ($r_0\ge\diam M$), (H3) holds with no further
condition; if $R$ is constant, (H3) holds if and only if $c_F>0$, i.e.\ iff $F$ does
not identify any pair at distance $\ge r_0$.
\end{itemize}
\end{lemma}

\begin{proof}
(a) Suppose $N(x)=N(y)$ with $0<d:=d_M(x,y)\le r_0$. By Lemma~\ref{lem:align} and (H2),
$\lVert F(x)-F(y)\rVert=\lvert R(x)-R(y)\rvert\le\Lambda_0 d$; by (H1),
$\lVert F(x)-F(y)\rVert\ge d/L_0$. Hence $d/L_0\le\Lambda_0 d$, contradicting
$\Lambda_0<1/L_0$.

(b) Suppose $N(x)=N(y)$ with $d\ge r_0$. By Lemma~\ref{lem:align},
$\lVert F(x)-F(y)\rVert=\lvert R(x)-R(y)\rvert\le\delta$ (the two bounds
$\rho^+-\rho^-$ and $\Lambda_0 d\le\Lambda_0\diam M$ both apply). But
$\lVert F(x)-F(y)\rVert\ge c_F>\delta$: contradiction, so no far pair aligns; together
with (a), $N$ is injective. For the quantitative bound, apply Paper~I's normalization
identity to a far pair:
\[
\lVert N(x)-N(y)\rVert^2
=\frac{\lVert F(x)-F(y)\rVert^2-\bigl(R(x)-R(y)\bigr)^2}{R(x)R(y)}
\ \ge\ \frac{c_F^2-\delta^2}{(\rho^+)^2},
\]
using $\lVert F(x)-F(y)\rVert\ge c_F$, $\lvert R(x)-R(y)\rvert\le\delta$, and
$R(x)R(y)\le(\rho^+)^2$.

(c) If $r_0\ge\diam M$, every pair falls under case (a). If $R$ is constant then
$\rho^+=\rho^-$ and $\delta=0$, so (b) applies whenever $c_F>0$; conversely if $c_F=0$
some far pair has $F(x)=F(y)$, hence $N(x)=N(y)$, and (H3) fails. Compactness makes
$c_F$ an attained minimum, so $c_F>0$ is equivalent to ``$F$ identifies no pair at
distance $\ge r_0$''.
\end{proof}

\begin{remark}
Consequences for Paper~I. (i) The global corollary (Cor.\ ``Global conditional form'')
can replace hypothesis (H3) by the checkable pinching condition $c_F>\delta$, and its
nonconstructive constant $c_0$ becomes the explicit
$\sqrt{c_F^2-\delta^2}/\rho^+$, making $A_g$ fully quantitative. (ii) Portegies-type
mode counts give injectivity of $F$, hence $c_F>0$ by compactness; the remaining
condition is the pinching $c_F>\delta$. (iii) The homogeneous-space remark of Paper~I
is the case $\rho^+=\rho^-$ of (c). (iv) A distance-resolved refinement holds with the
same proof: alignment at distance $d$ is impossible whenever
$\min\{\lVert F(x)-F(y)\rVert:d_M(x,y)=d\}>\min(\rho^+-\rho^-,\Lambda_0 d)$.
\end{remark}

\section{Atomlessness is automatic}

Paper~I's rank-limit results (the flat-torus proposition and the Green-kernel rank-limit
theorem) assume that the laws of $d_M(X,Y)$ and $G(X,Y)$ are atomless. Both assumptions
can be deleted.

We use twice the standard level-set lemma: \emph{if $h\in C^1$ on an open subset of
$\mathbb R^d$ (or a manifold chart), then $\nabla h=0$ almost everywhere on
$\{h=c\}$} (see \S5). Iterating: if $f\in C^2$ and $E=\{f=c\}$ has positive measure,
then a.e.\ on $E$ we have $\nabla f=0$; since $E\subseteq\{\partial_i f=0\}$ up to a
null set for each $i$, a second application (to $h=\partial_i f$) gives
$\nabla\partial_i f=0$ a.e.\ on $E$; hence the coordinate Hessian vanishes a.e.\ on
$E$, and therefore so does the Laplace--Beltrami operator,
$\Delta_g f=g^{ij}(\partial_i\partial_j f-\Gamma^k_{ij}\partial_k f)=0$ a.e.\ on $E$.

\begin{lemma}[Automatic atomlessness]\label{lem:atomless}
Let $M$ be a closed smooth Riemannian manifold (any dimension $d\ge1$, volume
normalized to $1$), $G$ its zero-mean Green kernel, and let $(X,Y)$ have a law with
bounded density with respect to $\vol\otimes\vol$. Then the laws of $d_M(X,Y)$ and of
$G(X,Y)$ are atomless.
\end{lemma}

\begin{proof}
By the bounded-density assumption it suffices to show that for each fixed $x$ the sets
$\{y:d_M(x,y)=t\}$ and $\{y:G(x,y)=c\}$ are $\vol$-null; Fubini then gives
$\mathbb P(d_M(X,Y)=t)=\mathbb P(G(X,Y)=c)=0$.

\emph{Distance.} Fix $x$. The function $d_x:=d_M(x,\cdot)$ is $1$-Lipschitz, and on
$M\setminus(\{x\}\cup\mathrm{Cut}(x))$ it is smooth with $\lvert\nabla d_x\rvert=1$
(the eikonal identity); the cut locus $\mathrm{Cut}(x)$ is closed and $\vol$-null
(imported, \S5). By the level-set lemma applied to the Lipschitz function $d_x$,
$\nabla d_x=0$ a.e.\ on $\{d_x=t\}$. Off the null set $\{x\}\cup\mathrm{Cut}(x)$ this
contradicts $\lvert\nabla d_x\rvert=1$, so $\vol\{d_x=t\}=0$.

\emph{Green kernel.} Fix $y$. By elliptic regularity $G(\cdot,y)$ is smooth on
$M\setminus\{y\}$ and satisfies, with the positive Laplacian and unit volume,
$\Delta\,G(\cdot,y)=-1$ there. If $\{x:G(x,y)=c\}$ had positive measure, the iterated
level-set lemma above would give $\Delta G(\cdot,y)=0$ on a positive-measure subset,
contradicting $\Delta G(\cdot,y)\equiv-1$. So each slice is null.
\end{proof}

\begin{remark}
(i) Consequently the flat-torus rank proposition and the Green-kernel rank-limit
theorem of Paper~I hold under the single hypothesis of a bounded pair density: their
atomlessness assumptions are theorems, not hypotheses. (ii) The argument does not
apply to the \emph{truncated} kernels $G_\Lambda$ (there
$\Delta_x G_\Lambda(\cdot,y)=E_\Lambda(\cdot,y)$ is not a nonzero constant), but the
limit-law atomlessness is all those proofs use. (iii) The argument is
dimension-independent; it does not use $d\le3$.
\end{remark}

\section{Green-rank positivity is an open condition}

\begin{theorem}[Perturbative openness of Green-rank positivity]\label{thm:open}
Let $M_0$ be a closed smooth manifold of dimension $d\le3$ and let $g_0$ be a smooth
Riemannian metric on $M_0$, volume normalized to $1$. Then the map
$g\mapsto\rho_G(g):=\rho_S\bigl(d_g(X,Y),-G_g(X,Y)\bigr)$ (with $(X,Y)$ i.i.d.\ from
normalized $\vol_g$) is continuous at $g_0$ with respect to $C^2$ convergence of
metrics, and likewise for the Kendall version $\tau_G$. Consequently:
\begin{itemize}
\item[(a)] if $\rho_G(g_0)>0$, then $\rho_G>0$ on a $C^2$-neighborhood of $g_0$;
\item[(b)] in particular, every two-point homogeneous metric of dimension $\le3$
($S^2,S^3,\mathbb{RP}^2,\mathbb{RP}^3$; and every metric on $S^1$) has a
$C^2$-neighborhood on which $\rho_G>0$;
\item[(c)] combined with the Green-kernel rank-limit theorem of Paper~I (whose
atomlessness hypotheses are deleted by Lemma~\ref{lem:atomless}), for every $g$ in
such a neighborhood there exist $\Lambda_0(g)<\infty$ and $\rho_0(g)>0$ such that the
truncated commute angular ranking satisfies
$\rho_S\bigl(d_g(X,Y),\lVert N_\Lambda(X)-N_\Lambda(Y)\rVert\bigr)\ge\rho_0(g)$ for
every multiplet-closed $\Lambda\ge\Lambda_0(g)$: the continuum rank form of Angular
Preservation, uniform in the mode count, holds on a $C^2$-open set of metrics.
\end{itemize}
\end{theorem}

\begin{proof}
Throughout, $g\to g_0$ in $C^2$; after a global rescaling (continuous in $g$) we may
keep $\vol_g(M)=1$. Write $\vol_0=\vol_{g_0}$ and
$\theta_g=d\vol_g/d\vol_0\to1$ in $C^1$.

\emph{Step 1: distances.} If $(1+\eta)^{-1}g_0\le g\le(1+\eta)g_0$ as quadratic forms
(which holds with $\eta=\eta(g)\to0$), then
$(1+\eta)^{-1/2}d_{g_0}\le d_g\le(1+\eta)^{1/2}d_{g_0}$, so $d_g\to d_{g_0}$ uniformly
on $M\times M$.

\emph{Step 2: uniform spectral tails.} The Dirichlet forms
$\mathcal E_g(u)=\int\lvert\nabla u\rvert_g^2\,d\vol_g$ and the $L^2$ norms are
uniformly comparable with constants $(1+\eta)^{\pm c(d)}$, so by min--max
$\mu_k(g)\ge(1+\eta)^{-c(d)}\mu_k(g_0)$ for all $k$. Since $d\le3$ gives
$\sum_k\mu_k(g_0)^{-2}<\infty$, the tails
$\sum_{k>K}\mu_k(g)^{-2}\le(1+\eta)^{2c}\sum_{k>K}\mu_k(g_0)^{-2}$ are small uniformly
over the neighborhood once $K$ is large.

\emph{Step 3: heads via resolvents.} Let $J_g:L^2(\vol_g)\to L^2(\vol_0)$,
$J_gu=\theta_g^{1/2}u$, a unitary; $A_g:=J_g\Delta_gJ_g^{-1}$ is self-adjoint on the
fixed space $L^2(\vol_0)$ with the same spectrum, and its quadratic form
$a_g(v)=\mathcal E_g(\theta_g^{-1/2}v)$ has coefficients converging uniformly to those
of $a_{g_0}$ (this uses $\theta_g\to1$ in $C^1$, hence $g\to g_0$ in $C^2$ suffices;
$C^1$ is what the computation needs). Uniformly comparable closed forms converging in
this sense yield norm resolvent convergence $A_g\to A_{g_0}$ (imported, \S5). Fix $K$
as in Step~2 with $\mu_K(g_0)<\mu_{K+1}(g_0)$, and a contour
$\Gamma\subset\mathbb C\setminus\operatorname{spec}(A_{g_0})$ enclosing
$\{\mu_1(g_0),\dots,\mu_K(g_0)\}$ but not $0$ nor $\mu_{K+1}(g_0)$. For $g$ near
$g_0$, $\Gamma$ still separates the same eigenvalue count (eigenvalue continuity from
Step~2's two-sided comparison), and
\[
\mathrm{Head}_g:=\frac{1}{2\pi i}\oint_\Gamma z^{-1}(z-A_g)^{-1}\,dz
=\sum_{k\le K}\mu_k(g)^{-1}\,P_k(g)
\]
(the residue of $z^{-1}(z-A_g)^{-1}$ at an eigenvalue $\mu$ is $\mu^{-1}$ times its
spectral projector; $0$ lies outside $\Gamma$, so the constant mode does not
contribute). Norm resolvent convergence gives
$\mathrm{Head}_g\to\mathrm{Head}_{g_0}$ in operator norm; both sides have rank $K$, so
the difference has rank $\le2K$ and
$\lVert T\rVert_{\mathrm{HS}}\le\sqrt{\operatorname{rank}T}\,\lVert T\rVert_{\mathrm{op}}$
upgrades this to Hilbert--Schmidt convergence, i.e.\ $L^2(M\times M)$ convergence of
the head kernels. Undoing the unitary multiplies kernels by
$\theta_g^{-1/2}(x)\theta_g^{-1/2}(y)\to1$ uniformly. Combining with Step~2,
$G_g\to G_{g_0}$ in $L^2(M\times M,\vol_0\otimes\vol_0)$.

\emph{Step 4: joint weak convergence of the pair statistics.} For bounded continuous
$h$,
\[
\mathbb E_g\,h\bigl(d_g(X,Y),G_g(X,Y)\bigr)
=\int h\bigl(d_g,G_g\bigr)\,\theta_g(x)\theta_g(y)\,d\vol_0\,d\vol_0 .
\]
The densities $\to1$ uniformly; $d_g\to d_{g_0}$ uniformly; $G_g\to G_{g_0}$ in
$L^2(\vol_0^{\otimes2})$, hence in measure; so
$h(d_g,G_g)\to h(d_{g_0},G_{g_0})$ in measure and, being bounded, in $L^1(\vol_0^{\otimes2})$.
Thus
$\bigl(d_g(X,Y),G_g(X,Y)\bigr)\Rightarrow\bigl(d_{g_0}(X,Y),G_{g_0}(X,Y)\bigr)$.

\emph{Step 5: continuity of the rank functionals.} By Lemma~\ref{lem:atomless} the
marginal laws of $d_g(X,Y)$ and $G_g(X,Y)$ are atomless for \emph{every} smooth $g$,
so $\rho_G(g)$ is the honest grade correlation
$\rho_G(g)=12\,\mathbb E\bigl[F_{d_g}(d_g)\,F_{-G_g}(-G_g)\bigr]-3$. Marginal weak
convergence to atomless limits gives uniform convergence of the CDFs (P\'olya);
replacing $F_{d_g},F_{-G_g}$ by $F_{d_{g_0}},F_{-G_{g_0}}$ costs $o(1)$, and the
remaining expectation converges by Step~4 since
$(u,v)\mapsto F_{d_{g_0}}(u)F_{-G_{g_0}}(-v)$ is bounded and continuous. Hence
$\rho_G(g)\to\rho_G(g_0)$. The Kendall functional, written on two independent pairs,
converges by the same weak-convergence-plus-null-discontinuity argument used in
Paper~I's circle theorem.

\emph{Conclusions.} (a) is immediate. (b): two-point homogeneity gives
$\rho_G(g_0)=1$ by Paper~I's corollary; on $S^1$ every smooth metric is isometric to a
round circle (arc-length parametrization), so $\rho_G\equiv1$ there outright. (c)
combines (a) with the Green-kernel rank-limit theorem, whose only remaining hypothesis
(bounded pair density) holds for i.i.d.\ volume sampling.
\end{proof}

\begin{remark}
The neighborhood in Theorem~\ref{thm:open} is not effective: Step~3 uses qualitative
norm resolvent convergence. An effective radius would follow from quantitative
resolvent perturbation bounds in terms of $\lVert g-g_0\rVert_{C^1}$ and the spectral
gap at $\mu_K(g_0)$; this is routine but unpleasant, and is left as the quantitative
frontier.
\end{remark}

\section{An anti-concentration transfer theorem (graph level)}

The conditional theorem of Paper~I transfers the \emph{metric} through (H1)--(H2). For
the \emph{rank} clause none of that is needed: calibration plus anti-concentration
suffice, with no resolution cutoff $\theta_n$ and no bi-Lipschitz hypothesis.

Fix a multiplet-closed cutoff $\Lambda$ and write
$V(x,y)=\lVert N_\Lambda(x)-N_\Lambda(y)\rVert$ for the continuum angular distance and
$D(x,y)=d_M(x,y)$. For an i.i.d.\ sample $x_1,\dots,x_n$ from a bounded, bounded-below
density, let $\widehat V_{ij}$ be the graph angular distances and suppose the
calibration event of Paper~I's transfer lemma holds:
$\max_i\lVert\hat q_i-N_\Lambda(x_i)\rVert\le2E_n/\rho^-_\Lambda$, hence
\begin{equation}\label{eq:cal}
\lvert\widehat V_{ij}-V(x_i,x_j)\rvert\ \le\ \bar\varepsilon_n:=\frac{4E_n}{\rho^-_\Lambda}
\qquad\text{for \emph{every} pair } i\ne j
\end{equation}
(the per-node bound holds at every node; no distance restriction enters).

\smallskip
\noindent\textbf{(AC) Anti-concentration.} $\displaystyle
\kappa_\Lambda:=\sup_{x\in M}\ \bigl\lVert\text{density of } V(x,Y)\bigr\rVert_\infty<\infty$,
where $Y$ is drawn from the sampling density.

\begin{theorem}[Rank transfer without bi-Lipschitz hypotheses]\label{thm:ac-transfer}
Let $d\le3$ and suppose $\rho_G(M)>0$ (e.g.\ $M$ in the open set of
Theorem~\ref{thm:open}), so that by the Green-kernel rank-limit theorem there are
$\Lambda_0$, $\rho_\star>0$, $\tau_\star>0$ with
$\rho_S(D,V)\ge\rho_\star$ and $\tau_K(D,V)\ge\tau_\star$ for every multiplet-closed
$\Lambda\ge\Lambda_0$. Fix such a $\Lambda$ and assume (AC) and the calibration
\eqref{eq:cal} with $E_n\to0$. Then, with probability tending to one, the
\emph{empirical} rank correlations of the graph angular distances against the manifold
distances over all sampled pairs satisfy
\[
\widehat\tau_K\bigl(\widehat V,D\bigr)\ \ge\ \tau_\star-C\,\kappa_\Lambda\bar\varepsilon_n-o(1),
\qquad
\widehat\rho_S\bigl(\widehat V,D\bigr)\ \ge\ \rho_\star-C\,\kappa_\Lambda\bar\varepsilon_n-o(1),
\]
with an absolute constant $C$; in particular both are eventually bounded below by
$\tau_\star/2$ and $\rho_\star/2$ respectively.
\end{theorem}

\begin{remark}[(AC) is to be read as the modulus form (AC$'$)]\label{rem:ac-super}
The bounded-density hypothesis (AC) as displayed is \emph{false} in general---a nondegenerate
saddle of $V(x,\cdot)$ forces a van~Hove logarithmic divergence of $\kappa_\Lambda$
(\S\ref{sec:ac}, guaranteed on any closed surface of Euler characteristic $\le0$). The
proof, however, uses only the weaker \emph{modulus} form
\[
\text{(AC$'$)}\qquad
\sup_{x\in M}\sup_{t}\ \mathbb P_Y\bigl(\lvert V(x,Y)-t\rvert\le\varepsilon\bigr)\ \le\ \omega(\varepsilon),
\quad \omega(0^+)=0,
\]
and the conclusion holds verbatim with $\kappa_\Lambda\bar\varepsilon_n$ replaced by
$\omega(2\bar\varepsilon_n)$, hence under $\omega(2\bar\varepsilon_n)\to0$. \S\ref{sec:ac}
establishes this substitution and proves (AC$'$) for every real-analytic metric; throughout,
(AC) is to be understood as (AC$'$).
\end{remark}

\begin{proof}
Write $N_p=\binom n2$ for the number of pairs, $V_P$ for the continuum value of pair
$P$ and $\widehat V_P$ for the graph value; ties are handled by the usual $\sgn$
conventions and occur with probability zero on the continuum side
(Lemma~\ref{lem:atomless} applies to $V$ through its strictly monotone relation to the
kernel $C_\Lambda$? --- not needed: we only use that flips are counted, see below).

\emph{Step 1: perturbation flips few comparisons.} For a pair-of-pairs $\{P,Q\}$, the
comparison $\sgn(\widehat V_P-\widehat V_Q)$ can differ from $\sgn(V_P-V_Q)$ only if
$\lvert V_P-V_Q\rvert\le2\bar\varepsilon_n$, by \eqref{eq:cal}. Let $f_n$ be the
fraction of pairs-of-pairs with $\lvert V_P-V_Q\rvert\le2\bar\varepsilon_n$. Then
\[
\bigl\lvert\widehat\tau_K(\widehat V,D)-\widehat\tau_K(V,D)\bigr\rvert\le2f_n .
\]
For Spearman: the $V$-ranking and the $\widehat V$-ranking of the $N_p$ pairs differ
by a permutation whose Kendall distance is at most $f_n\binom{N_p}2$; by the
Diaconis--Graham inequality the footrule (total rank displacement) is at most twice
the Kendall distance, and since one unit of rank displacement changes
$\sum_i d_i^2$ in the Spearman formula
$\widehat\rho_S=1-6\sum d_i^2/(N_p^3-N_p)$ by at most $2N_p$, we get
\[
\bigl\lvert\widehat\rho_S(\widehat V,D)-\widehat\rho_S(V,D)\bigr\rvert
\ \le\ \frac{6\cdot 2N_p\cdot 2 f_n\binom{N_p}2}{N_p^3-N_p}\ \le\ 12f_n+O(1/N_p).
\]

\emph{Step 2: the flip fraction is small in expectation.} Take a pair-of-pairs
$\{P,Q\}$. If $P,Q$ are vertex-disjoint, condition on $V_Q$: by (AC), integrating the
conditional density of $V_P$ given one endpoint (then averaging over it) bounds the
conditional density of $V_P$ by $\kappa_\Lambda$, so
$\mathbb P(\lvert V_P-V_Q\rvert\le2\bar\varepsilon_n)\le4\kappa_\Lambda\bar\varepsilon_n$.
If $P=(x_i,x_j)$ and $Q=(x_i,x_k)$ share the anchor $x_i$, condition on $x_i$ and
$V_Q$: given $x_i$, the value $V_P=V(x_i,x_j)$ with $x_j$ independent has density
$\le\kappa_\Lambda$ by (AC) (this is exactly why (AC) is stated uniformly in the
anchor), so the same bound holds. Hence
$\mathbb E f_n\le4\kappa_\Lambda\bar\varepsilon_n$, and by Markov
$f_n\le\kappa_\Lambda\bar\varepsilon_n^{1/2}\cdot$(const), say, with probability
tending to one; any rate with $f_n=O_P(\kappa_\Lambda\bar\varepsilon_n)$ suffices
below.

\emph{Step 3: empirical converges to population.} $\widehat\tau_K(V,D)$ is (up to
$O(1/n)$) a U-statistic of order four in the i.i.d.\ sample, so
$\widehat\tau_K(V,D)\to\tau_K(D,V)$ in probability (imported, \S5); the analogous
consistency of the sample Spearman coefficient for continuous population laws is
classical (\S5), and the continuum laws here are atomless
(Lemma~\ref{lem:atomless} for $D$; for $V$, the population functionals are those of
the Green-limit theorem, which are well-defined grade correlations there).

\emph{Assemble.} On the intersection of the calibration event, the Step-2 event, and
the Step-3 convergence,
$\widehat\tau_K(\widehat V,D)\ge\tau_K(D,V)-2f_n-o_P(1)
\ge\tau_\star-C\kappa_\Lambda\bar\varepsilon_n-o_P(1)$, and identically for
$\widehat\rho_S$ with the constant from Step~1.
\end{proof}

\begin{remark}[The growing-window condition, rederived]\label{rem:EA}
At a growing cutoff, the angular distances concentrate: $V=\sqrt2+O(A_\Lambda^{-1})$
with the pair-dependence of size $\asymp A_\Lambda^{-1}$ (Paper~I). Concentration
forces $\kappa_\Lambda\gtrsim A_\Lambda$, and the natural scale-aware form of (AC) is
the matching upper bound $\kappa_\Lambda\le C\,A_\Lambda$. Since
$\rho^-_\Lambda\asymp\sqrt{A_\Lambda}$, the error term in
Theorem~\ref{thm:ac-transfer} scales as
\[
\kappa_\Lambda\,\bar\varepsilon_n\ \asymp\ A_\Lambda\cdot\frac{E_n}{\sqrt{A_\Lambda}}
\ =\ E_n\sqrt{A_\Lambda},
\]
so the transfer condition is $E_n\sqrt{A_m}\to0$ --- \emph{exactly} the sharpened
growing-window condition derived in Paper~I by the direct kernel comparison. That two
independent routes (kernel rescaling; anti-concentration counting) produce the same
condition is evidence that $E_n\sqrt{A_m}\to0$ is the correct calibration threshold,
not an artifact of either proof.
\end{remark}

\begin{remark}[Status of (AC)]
(AC) at fixed $\Lambda$ asserts a bounded density for the pushforward of the sampling
measure under $y\mapsto\lVert N_\Lambda(x)-N_\Lambda(y)\rVert$, uniformly in $x$. By
the coarea formula this reduces to an integrated gradient lower bound
$\int_{\{V(x,\cdot)=t\}}\lvert\nabla_yV(x,y)\rvert^{-1}$ being bounded, i.e.\ to
$\nabla_yV$ not degenerating on large portions of level sets --- eigenfunction-level
technology (nodal/critical set estimates) is the natural tool, and a proof is a
self-contained open problem. Two mitigations: (i) (AC) is directly measurable on any
substrate (estimate the density of sampled $V$-values; the pre-registered pipeline can
emit this per run), turning the hypothesis into a per-instance certificate; (ii)
Theorem~\ref{thm:ac-transfer} survives with (AC) replaced by any bound
$\mathbb P(\lvert V_P-V_Q\rvert\le\epsilon)\le\omega(\epsilon)$ with
$\omega(2\bar\varepsilon_n)\to0$ --- bounded density is the simplest sufficient form,
not the needed one.
\end{remark}

\begin{remark}[Graph geodesics on the $D$-side]
The theorem compares against manifold distances $D=d_M$. Replacing $D$ by graph-hop
geodesics adds a $D$-side perturbation, handled by the identical Step-1/Step-2
argument once a uniform graph-geodesic consistency bound is available (standard for
$\varepsilon$-graphs above the connectivity scale); the required anti-concentration on
the $D$-side is automatic, since the pair-distance law on a closed manifold has a
bounded density.
\end{remark}



\section{The anti-concentration hypothesis (AC$'$)}\label{sec:ac}
\input{_sec_ac}

\section{A non-radiality calculus for Green-rank positivity}\label{sec:positivity}
\input{_sec_master}

\section{A geometric bound: the heat-kernel shell-variance}\label{sec:l2}

\noindent\textbf{Purpose.} The distributional master inequality (companion note
``green-rank-positivity'', Theorem ``Distributional master inequality'') gives
$\rhoG(M)\ge1-18(\phi_\Psi^2V_\Psi)^{1/3}$, reducing Green-rank positivity to a bound on the
$L^2$ non-radiality $V_\Psi=\E[(-G-\Psi\circ\dM)^2]$. That note left $V$ as a quantity to be
``bounded by heat-kernel anisotropy'' without proof. This note supplies the reduction. Its
three rigorous outputs: (i) the optimal $V$ is exactly the \emph{within-distance-shell
variance of the Green kernel}, $V_{\min}=\E[\Var(G\mid\dM)]$; (ii) via the heat
representation it equals the integrated \emph{shell-variance of the heat kernel}, which
\emph{vanishes identically on two-point homogeneous spaces} and is dominated by a heat-trace
envelope whose integral converges \emph{exactly for $d\le3$}; and (iii) an explicit
large-time spectral-gap tail plus a quantitative perturbative bound
$V_{\min}\le C(g_0)\lVert g-g_0\rVert_{C^2}^2$ near any two-point homogeneous $g_0$. The one
non-explicit constant---the small-time curvature-anisotropy amplitude---is isolated and
named. Notation follows the companion notes; $(X,Y)$ are i.i.d.\ from the unit-volume
measure and all $L^2$ norms are over $M\times M$ with that measure.

\section{The optimal profile is the distance-conditional mean}

\begin{proposition}[Conditional-variance characterization]\label{prop:condvar}
Over all (Borel) radial profiles $\Psi$, $V_\Psi=\E[(-G(X,Y)-\Psi(\dM(X,Y)))^2]$ is minimized
by the distance-conditional mean $\Psi^\star(r)=-\E[G(X,Y)\mid\dM(X,Y)=r]$, and
\[
V_{\min}:=V_{\Psi^\star}=\E\bigl[\Var(G(X,Y)\mid\dM(X,Y))\bigr].
\]
Near the diagonal $\Psi^\star$ is strictly increasing with $\phi_{\Psi^\star}<\infty$; its
global admissibility is condition~\textup{(A$^\star$)} below.
\end{proposition}

\begin{proof}
For fixed $r$ the map $c\mapsto\E[(-G-c)^2\mid\dM=r]$ is minimized at $c=\E[-G\mid\dM=r]$
with minimal value $\Var(-G\mid\dM=r)=\Var(G\mid\dM=r)$; averaging over the law of $\dM(X,Y)$
gives the claim (the standard $L^2$-projection onto $\sigma(\dM)$). For the near-diagonal
regularity: writing $G=\gamma_d\circ\dM+h$ (parametrix lemma, companion note),
$\E[G\mid\dM=r]=\gamma_d(r)+\E[h\mid\dM=r]$ with $h$ bounded, so
$\Psi^\star=-\gamma_d-\E[h\mid\dM=\cdot]$ has the strictly \emph{increasing} singular part
$-\gamma_d$ (as $\gamma_d$ decreases in $r$) dominating a bounded correction; the
$\phi_{\Psi^\star}<\infty$ computation is then the parametrix change-of-variables verbatim.
\end{proof}

\begin{remark}[The admissibility condition (A$^\star$)]\label{rem:astar}
The bound below controls the \emph{unconstrained} minimizer $V_{\min}$, whose profile
$\Psi^\star$ must itself be a monotone profile with finite $\phi$ before $V_{\min}$ may enter
the distributional master inequality. We isolate this as
\smallskip
\noindent\textbf{(A$^\star$)}\quad \emph{$r\mapsto\E[G\mid\dM=r]$ is strictly decreasing on
$(0,\diam M)$, with $\phi_{\Psi^\star}<\infty$.}
\smallskip
By the Proposition, (A$^\star$) holds near the diagonal unconditionally; and it holds on a
$C^2$-neighborhood of any two-point homogeneous metric, where $\E[G\mid\dM]=g(\dM)$ is
strictly decreasing (Paper~I) and strict monotonicity plus $\phi<\infty$ persist under small
$C^2$ perturbation. It is a checkable, measurable condition (evaluate the shell-mean of $G$).
Away from these regimes it is a hypothesis; where it fails one falls back on the fixed
parametrix profile at the cost of losing the vanishing on two-point homogeneous spaces.
\end{remark}

\section{Heat representation: the shell-variance of the heat kernel}

Let $p_t(x,y)=\sum_{k\ge0}e^{-\mu_k t}\varphi_k(x)\varphi_k(y)$ be the heat kernel
($\mu_0=0$, $\varphi_0\equiv1$), so $G(x,y)=\int_0^\infty(p_t(x,y)-1)\,dt$. Write the
shell-centering $P_{\mathrm{rad}}F(x,y)=\E[F\mid\dM]=F$'s conditional mean given distance,
and $\delta_t=p_t-P_{\mathrm{rad}}p_t$ (the part of the time-$t$ heat kernel that is
\emph{not} a function of distance). Define the \emph{time-$t$ heat-kernel non-radiality}
\[
\mathcal N(t):=\lVert\delta_t\rVert_{L^2}=\bigl(\E[\Var(p_t(X,Y)\mid\dM(X,Y))]\bigr)^{1/2}.
\]

\begin{lemma}[Reduction to the heat-kernel non-radiality]\label{lem:reduce}
\[
V_{\min}^{1/2}\ \le\ \int_0^\infty\mathcal N(t)\,dt,
\qquad\text{and}\qquad
\mathcal N(t)\ \le\ \bigl(Z(2t)-1\bigr)^{1/2},
\]
where $Z(s)=\Tr e^{-s\Delta}=\sum_{k\ge0}e^{-\mu_k s}$ is the heat trace. Consequently
$\mathcal N(t)\equiv0$ when $M$ is two-point homogeneous (there $p_t$ is a function of
distance for every $t$), and the envelope integral $\int_0^\infty(Z(2t)-1)^{1/2}\,dt$
converges \emph{iff} $d\le3$.
\end{lemma}

\begin{proof}
The optimal remainder is $R^\star=-G-\Psi^\star\circ\dM=-(G-\E[G\mid\dM])
=-\int_0^\infty(p_t-P_{\mathrm{rad}}p_t)\,dt=-\int_0^\infty\delta_t\,dt$, the interchange of
$\int dt$ and the (bounded, linear) conditional expectation being justified by absolute
integrability of $p_t-1$ in $L^2$ (below). Minkowski's integral inequality gives
$V_{\min}^{1/2}=\lVert R^\star\rVert_2\le\int_0^\infty\lVert\delta_t\rVert_2\,dt
=\int_0^\infty\mathcal N(t)\,dt$.

For the envelope: shell-centering is an $L^2$-orthogonal projection, hence a contraction, so
by the law of total variance $\mathcal N(t)^2=\E[\Var(p_t\mid\dM)]\le\Var(p_t(X,Y))
=\E[p_t^2]-1=\lVert p_t\rVert_2^2-1=\sum_{k\ge0}e^{-2\mu_k t}-1=Z(2t)-1$ (using
$\E[p_t(X,Y)]=\int\!\!\int p_t=1$ and Parseval). Two-point homogeneity makes the invariant
kernel $p_t$ a function of geodesic distance, so $\delta_t\equiv0$ and $\mathcal N\equiv0$.
Finally $Z(2t)-1=\sum_{k\ge1}e^{-2\mu_k t}$: as $t\to0$, $Z(2t)\sim(8\pi t)^{-d/2}$ (Weyl
heat-trace asymptotics, $\vol=1$), so $(Z(2t)-1)^{1/2}\sim c\,t^{-d/4}$, integrable at $0$
iff $d/4<1$, i.e.\ $d\le3$; as $t\to\infty$, $Z(2t)-1\le C e^{-2\mu_1 t}$, integrable. Hence
the envelope integral converges iff $d\le3$, and it dominates $\int\mathcal N$, which is
therefore finite for $d\le3$ (justifying the interchange above).
\end{proof}

The critical dimension $d=4$ reappears exactly: the shell-variance envelope ceases to be
integrable, matching the Hilbert--Schmidt threshold and the fractional-filter remedy.

\section{The explicit large-time tail and the small-time anisotropy}

Split $\int_0^\infty\mathcal N=\int_0^{t_0}\mathcal N+\int_{t_0}^\infty\mathcal N$. The tail is
fully explicit; the head is reduced to a named curvature-anisotropy amplitude.

\begin{lemma}[Explicit spectral-gap tail]\label{lem:tail}
For any $t_0>0$,
$\displaystyle\int_{t_0}^\infty\mathcal N(t)\,dt\le\frac{(Z(2t_0)-1)^{1/2}}{\mu_1}$.
Under $\Ric\ge-(d-1)K$ and $\diam\le D$ the spectral gap obeys $\mu_1\ge\mu_1(K,D)>0$
(Li--Yau/Zhong--Yang), and $Z(2t_0)-1\le C(d,K,D)$ for $t_0\ge1$; so the tail is bounded by
an explicit function of $(d,K,D,t_0)$ decaying like $e^{-\mu_1 t_0}/\mu_1$.
\end{lemma}

\begin{proof}
For $t\ge t_0$, $Z(2t)-1=\sum_{k\ge1}e^{-2\mu_k t}\le e^{-2\mu_1(t-t_0)}\sum_{k\ge1}e^{-2\mu_k t_0}
=e^{-2\mu_1(t-t_0)}(Z(2t_0)-1)$, so $\mathcal N(t)\le(Z(2t_0)-1)^{1/2}e^{-\mu_1(t-t_0)}$ by
Lemma~\ref{lem:reduce}; integrating $\int_{t_0}^\infty e^{-\mu_1(t-t_0)}\,dt=1/\mu_1$ gives the
bound. The gap and heat-trace bounds are standard under the stated curvature/diameter
hypotheses.
\end{proof}

\begin{definition}[Curvature-anisotropy amplitude]\label{def:aniso}
Let $\mathcal A:=\sup_{0<t\le t_0}t^{d/4}\mathcal N(t)$, the supremal shell-standard-deviation
of the heat kernel on the head interval, normalized by the diagonal Weyl scale $t^{-d/4}$.
\end{definition}

\begin{lemma}[Head bound]\label{lem:head}
For $d\le3$, $\displaystyle\int_0^{t_0}\mathcal N(t)\,dt\le\mathcal A\int_0^{t_0}t^{-d/4}\,dt
=\frac{4\,\mathcal A}{4-d}\,t_0^{(4-d)/4}$. Here $\mathcal A<\infty$ always (since
$\mathcal N(t)\le(Z(2t)-1)^{1/2}$, so $t^{d/4}\mathcal N(t)$ is bounded as $t\to0$) and
$\mathcal A=0$ on two-point homogeneous spaces (there $\mathcal N\equiv0$). Its
\emph{small-time} part $\lim_{t\to0}t^{d/4}\mathcal N(t)$ is $O$ of the shell-oscillation of
the Minakshisundaram--Pleijel amplitude $u_0(x,y)=\det(d\exp)^{-1/2}$---the variation of the
Jacobi/curvature data along equal-length geodesics, $O(\sup\lVert\nabla\Ric\rVert$ and the
sectional-curvature variation$)$, vanishing when the curvature is isotropic and parallel (the
locally symmetric case). The \emph{full} sup $\mathcal A$ over $(0,t_0]$ is the manifold's
genuine heat-kernel non-radiality on that interval---generically $\Theta(1)$ for a
non-symmetric metric; making it curvature-explicit uniformly in $t\le t_0$ is the frontier
item, not claimed here.
\end{lemma}

\begin{proof}
The integral bound is Definition~\ref{def:aniso} and $\int_0^{t_0}t^{-d/4}dt=\frac{4}{4-d}
t_0^{(4-d)/4}$ (finite for $d<4$). For the amplitude: near the diagonal the MP expansion
$p_t(x,y)=(4\pi t)^{-d/2}e^{-\dM(x,y)^2/4t}\bigl(u_0(x,y)+u_1(x,y)t+\cdots\bigr)$ has a
Gaussian factor that is radial; the non-radial part of $p_t$ at leading order is
$(4\pi t)^{-d/2}e^{-\dM^2/4t}\bigl(u_0(x,y)-\E[u_0\mid\dM]\bigr)$, whose $L^2$-shell-variance
scaled by $t^{d/4}$ is $O(\text{shell-oscillation of }u_0)$; $u_0=\det(d\exp_x(y))^{-1/2}$ is
a function of the geometry along the $x$-to-$y$ geodesic, constant on distance shells exactly
when the geometry is isotropic and parallel, and its shell-oscillation is $O(d(x,y)^2)$ times
the curvature variation by the Jacobi-field (Van~Vleck--Morette) expansion
$u_0=1+\tfrac1{12}\Ric(\dot\gamma,\dot\gamma)r^2+O(r^3)$ (the volume density is
$\theta=\sqrt{\det g}=1-\tfrac16\Ric(\dot\gamma,\dot\gamma)r^2+O(r^3)$, and
$u_0=\theta^{-1/2}=(\det g)^{-1/4}$).
On a two-point homogeneous space $u_0$ is a function of distance, so $\mathcal A=0$; the
intermediate-time ($t\le t_0$) contribution is likewise $0$ there since $\mathcal N\equiv0$.
\end{proof}

\begin{theorem}[Geometric $L^2$ non-radiality bound]\label{thm:main}
Let $M$ be closed of dimension $d\le3$, $\vol(M)=1$, $\Ric\ge-(d-1)K$, $\diam\le D$,
injectivity radius $\ge i_0$. Then for every $t_0\ge1$,
\[
V_{\min}^{1/2}\ \le\ \frac{4\,\mathcal A}{4-d}\,t_0^{(4-d)/4}\ +\ \frac{(Z(2t_0)-1)^{1/2}}{\mu_1},
\]
with $\mu_1\ge\mu_1(K,D)>0$ and $Z(2t_0)-1\le C(d,K,D)$ explicit, and $\mathcal A$ the
curvature-anisotropy amplitude of Definition~\ref{def:aniso}. Consequently:
\begin{itemize}
\item[(i)] $V_{\min}=0$ on every two-point homogeneous space ($\mathcal A=0$ and
$\mathcal N\equiv0$), recovering $\rhoG=1$ there;
\item[(ii)] under \textup{(A$^\star$)} (Remark~\ref{rem:astar}), combined with the
distributional master inequality, $\rhoG(M)\ge1-18\bigl(\phi_{\Psi^\star}^2V_{\min}\bigr)^{1/3}$,
so $\rhoG(M)>0$ whenever the right side of the displayed bound is small enough that
$\phi_{\Psi^\star}^2V_{\min}<18^{-3}$---a class explicit in $(K,D,i_0,\mu_1)$ modulo the single
anisotropy amplitude $\mathcal A$.
\end{itemize}
\end{theorem}

\begin{proof}
Sum Lemmas~\ref{lem:head} and~\ref{lem:tail} into Lemma~\ref{lem:reduce}. (i) is
$\mathcal A=0$ and $\mathcal N\equiv0$; (ii) is the companion distributional inequality
applied to $\Psi^\star$, which is a monotone profile with $\phi_{\Psi^\star}<\infty$ exactly
under \textup{(A$^\star$)}.
\end{proof}

\begin{corollary}[Perturbative $\eta^{2/3}$ rate]\label{cor:pert}
Let $g_0$ be a two-point homogeneous metric ($\vol=1$) and $g$ a metric with
$\lVert g-g_0\rVert_{C^2}\le\eta$ small. Then $V_{\min}(g)\le C(g_0)\,\eta^2$, with $C(g_0)$
the standard parabolic-perturbation constant of $g_0$, and hence
$\rhoG(g)\ge1-18\,\bigl(\phi_0^2C(g_0)\bigr)^{1/3}\eta^{2/3}\to1$ as $\eta\to0$, quantifying
the perturbative openness theorem in the $L^2$ index.
\end{corollary}

\begin{proof}
By Remark~\ref{rem:astar}, (A$^\star$) holds for $g$ near $g_0$, so $V_{\min}(g)$ is
admissible. Since $g_0$ is two-point homogeneous, $\delta_t^{g_0}\equiv0$, so
$\mathcal N^g(t)=\lVert\delta_t^g\rVert_2=\lVert(I-P^g_{\mathrm{rad}})p_t^g
-(I-P^{g_0}_{\mathrm{rad}})p_t^{g_0}\rVert_2$; adding and subtracting,
\[
\mathcal N^g(t)\ \le\ \underbrace{\lVert p_t^g-p_t^{g_0}\rVert_2}_{\text{(I)}}
\ +\ \underbrace{\lVert(P^g_{\mathrm{rad}}-P^{g_0}_{\mathrm{rad}})p_t^{g_0}\rVert_2}_{\text{(II)}}.
\]
For (II): $\lVert d^g-d^{g_0}\rVert_\infty=O(\eta)$ (the metric is $C^0$-$O(\eta)$-close), so
the shell-centering conditional expectations differ by $O(\eta)$ applied to the radial-regular
$p_t^{g_0}$, giving (II)$\le C\eta\,(Z^{g_0}(2t)-1)^{1/2}$. For (I): the parabolic (Duhamel)
perturbation estimate---$p_t^g-p_t^{g_0}=-\int_0^t p^g_{t-s}(\Delta_g-\Delta_{g_0})p^{g_0}_s\,ds$
with $\Delta_g-\Delta_{g_0}$ of order two and coefficients $O(\lVert g-g_0\rVert_{C^1})$---gives
(I)$\le C\eta\,(Z^{g_0}(c t)-1)^{1/2}$ (imported fact~6). Both envelopes are $\asymp t^{-d/4}$
as $t\to0$ (integrable for $d\le3$) and decay like $e^{-c\mu_1(g_0)t}$ as $t\to\infty$; hence
$V_{\min}(g)^{1/2}=\int_0^\infty\mathcal N^g(t)\,dt\le C(g_0)^{1/2}\eta$. This uses only the
leading Duhamel bound---not smoothness of the higher heat coefficients $u_j$, which
$\lVert g-g_0\rVert_{C^2}$ does not control.
\end{proof}



\section{Past the critical dimension: the fractional filter}\label{sec:fractional}

\noindent\textbf{Purpose.} The rank form of Angular Preservation is settled below the
critical dimension $d=4$: the commute angular ranking converges uniformly in the mode
count to the Green-kernel ranking (Paper~I, Green-kernel rank-limit theorem), reducing the
question to Green-rank positivity. At $d=4$ the commute Green kernel
$G=\sum_k\mu_k^{-1}\varphi_k\otimes\varphi_k$ ceases to be Hilbert--Schmidt
($\sum_k\mu_k^{-2}=\infty$ by Weyl's law $\mu_k\asymp k^{2/d}$), so the limit step fails.
Paper~I flags the remedy---a stronger filter $\mu_k^{-s/2}$, $s>d/4$---and the first
question, monotonicity of the fractional Green kernel on two-point homogeneous spaces. This Part carries the entire chain up one level to the fractional filter and, using
heat-kernel subordination, proves the two-point homogeneous case \emph{in every dimension},
so that the reduction ``Angular Preservation $\Leftarrow$ fractional-Green-rank positivity''
holds for all $d$, with the geometric coefficient equal to $1$ on the rank-one symmetric
spaces regardless of dimension. Notation follows Paper~I and \S\ref{sec:positivity}.

\section{The fractional filter and Hilbert--Schmidt stability}

Fix $s>0$ and, at a multiplet-closed cutoff $\Lambda$, define the $s$-filtered embedding and
its kernel
\[
F_{s,\Lambda}=\bigl(\mu_k^{-s/2}\varphi_k\bigr)_{0<\mu_k\le\Lambda},\quad
K_{s,\Lambda}(x,y)=\!\!\sum_{0<\mu_k\le\Lambda}\!\!\frac{\varphi_k(x)\varphi_k(y)}{\mu_k^{s}},\quad
S_{s,\Lambda}(x)=K_{s,\Lambda}(x,x),\quad
N_{s,\Lambda}=\frac{F_{s,\Lambda}}{\sqrt{S_{s,\Lambda}}} .
\]
The commute filter is $s=1$. Write $K_s=\sum_{k\ge1}\mu_k^{-s}\varphi_k\otimes\varphi_k$ for
the full $s$-kernel (the kernel of $\Delta^{-s}$ on the zero-mean subspace) and
$\rho_{G,s}(M):=\rho_S\bigl(\dM(X,Y),-K_s(X,Y)\bigr)$ for the \emph{fractional Green-rank
coefficient}.

\begin{lemma}[Hilbert--Schmidt stability]\label{lem:hs}
$\lVert K_{s,\Lambda'}-K_{s,\Lambda}\rVert_{L^2(M\times M)}^2=\sum_{\Lambda<\mu_k\le\Lambda'}\mu_k^{-2s}$,
and by Weyl's law $\mu_k\asymp k^{2/d}$ this converges iff $s>d/4$. Hence for $s>d/4$,
$K_{s,\Lambda}\to K_s$ in $L^2(M\times M)$, and $K_s\in L^2$.
\end{lemma}

\begin{proof}
The $\varphi_k\otimes\varphi_k$ are orthonormal in $L^2(M\times M)$, giving the stated norm;
$\sum_k\mu_k^{-2s}\asymp\sum_k k^{-4s/d}$ converges iff $4s/d>1$.
\end{proof}

For $d\ge4$ any $s>d/4>1$ works; the minimal admissible filter grows with dimension.

\section{The fractional rank limit}

\begin{theorem}[Fractional Green-kernel rank limit]\label{thm:fraclimit}
Let $M$ be closed of dimension $d$, $\vol(M)=1$, and $s>d/4$. Let $(X,Y)$ have bounded
density with respect to $\vol\otimes\vol$, with the laws of $\dM(X,Y)$ and $K_s(X,Y)$
atomless. Then
\[
\rho_S\bigl(\dM(X,Y),\lVert N_{s,\Lambda}(X)-N_{s,\Lambda}(Y)\rVert\bigr)\ \longrightarrow\
\rho_S\bigl(\dM(X,Y),-K_s(X,Y)\bigr)=\rho_{G,s}(M),
\]
and likewise for Kendall's $\tau$. Consequently, if $\rho_{G,s}(M)>0$ there are
$\Lambda_0<\infty,\rho_0>0$ with the $s$-filtered angular ranking $\ge\rho_0$ for every
multiplet-closed $\Lambda\ge\Lambda_0$.
\end{theorem}

\begin{proof}
Identical in structure to the $s=1$ theorem (Paper~I, Green-kernel rank-limit), with
$\mu_k^{-1}$ replaced by $\mu_k^{-s}$ throughout. \emph{(1)} Hilbert--Schmidt limit is
Lemma~\ref{lem:hs}. \emph{(2)} The diagonal $S_{s,\Lambda}(x)=\sum_{\mu_k\le\Lambda}\mu_k^{-s}\varphi_k(x)^2$
is asymptotically constant uniformly in $x$: the uniform pointwise Weyl law
$\sum_{\mu_k\le L}\varphi_k(x)^2=c_dL^{d/2}+O(L^{(d-1)/2})$ (H\"ormander, uniform in $x$) and
partial summation give, \emph{for $d/4<s\le d/2$}, $S_{s,\Lambda}(x)=A_{s,\Lambda}+o(A_{s,\Lambda})$
uniformly in $x$, where $A_{s,\Lambda}$ is the leading term of $\int^\Lambda L^{-s}\,dc_dL^{d/2}$:
$\asymp\Lambda^{d/2-s}$ (diverges) for $s<d/2$, $\asymp\log\Lambda$ for $s=d/2$ (the uniform
Hörmander remainder $O(L^{(d-1)/2})$ contributes a relative $O(\Lambda^{-1/2})$, lower order).
For $s>d/2$ the diagonal instead \emph{converges} to a finite, generally non-constant $S_s(x)$
and does not rescale to a constant---this case is the Remark. \emph{(3)} Rescale $H_{s,\Lambda}=A_{s,\Lambda}C_{s,\Lambda}\to K_s$
in $L^2$, ranks unchanged by the positive constant. \emph{(4)} Bounded pair density carries
$L^2$-convergence to convergence in probability; atomlessness of the limit gives distribution-function
convergence and hence convergence of the population Spearman/Kendall coefficients.
\end{proof}

\begin{remark}[The convergent-diagonal regime $s>d/2$]
When $s>d/2$ the diagonal $S_{s,\Lambda}(x)\to S_s(x)=\sum_k\mu_k^{-s}\varphi_k(x)^2$, a
\emph{finite, generally non-constant} function of $x$ (the fractional analogue of a local
density). Step (2) then does not rescale to a constant; instead $C_{s,\Lambda}\to
K_s(x,y)/\sqrt{S_s(x)S_s(y)}$ pointwise and in $L^2$, and the rank limit is
$\rho_S(\dM,-K_s/\sqrt{S_sS_s})$. The normalization by $S_s$ is exactly the row-normalization
of the $s$-embedding; on a homogeneous space $S_s$ is constant and this reduces to
$\rho_{G,s}$. For the inhomogeneous case the operative coefficient is the
\emph{$S_s$-normalized} fractional Green-rank correlation, and Theorem~\ref{thm:fraclimit}
holds with $K_s$ replaced by $K_s/\sqrt{S_s(x)S_s(y)}$; all subsequent statements carry the
same replacement. We keep $K_s$ in the display for the homogeneous and $d/4<s\le d/2$ cases,
where $S_{s,\Lambda}$ genuinely rescales to a constant.
\end{remark}

\section{The fractional master inequality and the Riesz parametrix}

The master inequality is filter-agnostic: it used only that $-K$ is compared to a monotone
radial profile. Restating (\S\ref{sec:positivity}, Theorem~\ref{thm:master}): for any strictly increasing $\Psi$ with remainder
$\Rem_s=-K_s-\Psi\circ\dM$,
\[
\rho_{G,s}(M)\ \ge\ 1-3\,\mu_\Psi\!\bigl(2\lVert\Rem_s\rVert_\infty\bigr)\ \ge\
1-12\,\phi_\Psi\lVert\Rem_s\rVert_\infty,
\]
with $\rho_{G,s}=1$ when $\Rem_s\equiv0$. The near-diagonal remainder is again bounded, now
by the Riesz parametrix.

\begin{lemma}[Riesz parametrix]\label{lem:riesz}
Let $M$ be closed smooth, $\vol(M)=1$, and $s>(d-1)/2$ (a nonempty range for every $d$, and
automatically $>d/4$ since $d>2$). Near the diagonal
$K_s(x,y)=\gamma_{d,s}\,\dM(x,y)^{2s-d}+h_s(x,y)$ with $\gamma_{d,s}>0$ the universal Riesz
constant and $h_s$ bounded up to the diagonal; the leading profile
$r\mapsto\gamma_{d,s}r^{2s-d}$ (exponent $2s-d\in(-1,0)$, strictly decreasing) is radial with a
point-independent coefficient, and it is the \emph{only} singular part---the subleading
Hadamard corrections carry exponent $2s-d+2j\ge2s-d+2>1$ for $j\ge1$, hence are bounded even
though their coefficients are point-dependent. When $2s\ge d$, $K_s$ is bounded (continuous,
$2s>d$) or $-c\log\dM+$bounded ($2s=d$). In all cases $\Psi_s=-\gamma_{d,s}r^{2s-d}$ (resp.\
the bounded/log profile), extended strictly increasing on $[r_0,\diam M]$, has finite
$\lVert\Rem_{s}\rVert_\infty$ and finite $\phi_{\Psi_s}$.
\end{lemma}

\begin{proof}
$K_s=\Delta^{-s}$ on the zero-mean subspace is the Riesz potential of order $2s$; it is a
classical pseudodifferential operator of order $-2s$ with principal symbol
$\lvert\xi\rvert_g^{-2s}$, and its kernel has the Hadamard expansion
$K_s(x,y)=\sum_{j\ge0}c_j(x,y)\dM(x,y)^{2s-d+2j}+(\text{smooth})$ with $c_0=\gamma_{d,s}$ the
universal Euclidean Riesz constant (the metric is Euclidean to second order in normal
coordinates) and $c_j$ smooth for $j\ge1$. The restriction $s>(d-1)/2$ makes $2s-d>-1$, so
every $j\ge1$ term has exponent $2s-d+2j>1$ and is bounded up to the diagonal (with bounded,
point-dependent coefficient); thus only $c_0\dM^{2s-d}$ is singular, and
$h_s:=K_s-\gamma_{d,s}\dM^{2s-d}$ is bounded. Hence $\Rem_s=-K_s-\Psi_s\circ\dM=-h_s$ near the
diagonal, bounded; off the diagonal both $K_s$ and $\Psi_s\circ\dM$ are bounded. Finiteness of
$\phi_{\Psi_s}$ is the change-of-variables computation of \S\ref{sec:positivity} with exponent
$2s-d$ in place of $2-d$: $\Psi_s'\asymp r^{2s-d-1}$ and $f_W\asymp r^{d-1}$ give density
$\asymp r^{2(d-s)}$, which $\to0$ at $r\to0$ for $s<d$, and is bounded on $[r_0,\diam M]$.
\end{proof}

\begin{theorem}[Fractional near-radial positivity]\label{thm:fracnear}
Let $s>(d-1)/2$. With $\Psi_s$ the Riesz profile of Lemma~\ref{lem:riesz} (a \emph{single}
fixed profile, both $\phi_{\Psi_s}$ and $\lVert\Rem_s\rVert_\infty$ finite),
$\rho_{G,s}(M)\ge1-12\,\phi_{\Psi_s}\lVert\Rem_s\rVert_\infty$, so $\rho_{G,s}(M)>0$ whenever
$\phi_{\Psi_s}\lVert\Rem_s\rVert_\infty<1/12$; and by Theorem~\ref{thm:fraclimit} the
$s$-filtered angular ranking is then uniform-in-mode-count rank-faithful. The perturbative
openness argument (\S\ref{sec:reductions}) transfers with $\mu_k^{-1}\to\mu_k^{-s}$---its
spectral-tail step now requires $\sum_k\mu_k^{-2s}<\infty$ (i.e.\ $s>d/4$), the resolvent
contour weight becomes $z^{-s}$ (holomorphic off the negative axis, with $0$ excluded), and
atomlessness of the $K_s$-pair law is assumed as in Theorem~\ref{thm:fraclimit}---so
$\rho_{G,s}>0$ on a $C^2$-open set of metrics containing the exact cases of \S4.
\end{theorem}

\section{Two-point homogeneous spaces, all dimensions, via subordination}

The one place the fractional story is \emph{cleaner} than the commute case is the exact
computation: there is no local Laplace ODE for $\Delta^{-s}$, but subordination to the heat
semigroup supplies strict monotonicity in every dimension at once.

\begin{theorem}[Fractional Green-rank $=1$ on rank-one symmetric spaces]\label{thm:twopoint}
Let $M$ be a compact two-point homogeneous space (a sphere $S^d$, or a projective space
$\mathbb{RP}^d,\mathbb{CP}^d,\mathbb{HP}^d,\mathbb{OP}^2$) of any dimension, $\vol(M)=1$, and
let $s>0$. Then $K_s(x,y)=k_s(\dM(x,y))$ with $k_s$ strictly decreasing on
$(0,\diam M)$; hence $\rho_{G,s}(M)=1$, and the $S_s$-normalized coefficient of the Remark
is likewise $1$. In particular, for every $d$ and every admissible $s>d/4$, Angular
Preservation (fractional-filter rank form) holds on these spaces, uniformly in the mode
count.
\end{theorem}

\begin{proof}
By subordination, for $s>0$ and $\lambda>0$, $\lambda^{-s}=\frac1{\Gamma(s)}\int_0^\infty
t^{s-1}e^{-\lambda t}\,dt$; applied to $\Delta$ on the zero-mean subspace,
\[
K_s(x,y)=\frac1{\Gamma(s)}\int_0^\infty t^{s-1}\bigl(p_t(x,y)-1\bigr)\,dt,
\]
where $p_t$ is the heat kernel and the $-1$ removes the constant mode ($\vol M=1$). Two-point
homogeneity makes $p_t$ a function of geodesic distance, $p_t(x,y)=P_t(\dM(x,y))$, and on a
compact two-point homogeneous space $P_t$ is \emph{strictly decreasing} in the distance for
every $t>0$. This is the one genuinely analytic input; we justify it by the parabolic maximum
principle rather than by citation. Writing the radial heat equation
$\partial_tP=\partial_r^2P+m(r)\,\partial_rP$ with $m=A'/A$ the log-derivative of the
geodesic-sphere area density---on a rank-one symmetric space
$m(r)=\sum(\mathrm{mult})\cdot\tfrac{\alpha}{2}\cot\tfrac{\alpha r}{2}$, strictly
decreasing---the radial derivative $v:=\partial_rP_t$ solves
$\partial_tv=\partial_r^2v+m\,\partial_rv+m'(r)\,v$ with $v=0$ at the pole and the antipodal
focal set and $m'(r)<0$. Since $v<0$ for small $t$ (the kernel is a decreasing peak) and the
lower-order coefficient $m'$ is negative, the strong parabolic maximum principle keeps $v<0$
on the open interval $(0,\diam M)$ for all $t>0$ (individual zonal harmonics oscillate, but
the heat-weighted sum does not). Differentiating $K_s$ under the integral,
which is justified because $\partial_r P_t(r)$ is integrable against $t^{s-1}$ (small-$t$
Gaussian decay $\partial_rP_t=O(t^{-d/2-1}re^{-r^2/4t})$ times $t^{s-1}$ is integrable at
$0$ for $r>0$; large-$t$ exponential relaxation $\lvert\partial_rP_t\rvert\le Ce^{-\mu_1 t}$
times $t^{s-1}$ is integrable at $\infty$),
\[
k_s'(r)=\frac1{\Gamma(s)}\int_0^\infty t^{s-1}\,\partial_r P_t(r)\,dt\ <\ 0
\qquad(0<r<\diam M),
\]
since $t^{s-1}>0$ and $\partial_rP_t(r)<0$ for every $t>0$. Thus $-K_s$ is a strictly
increasing function of $\dM$, the pair $(\dM,-K_s)$ is comonotone, and $\rho_{G,s}=1$.
Normalizing by the constant diagonal $S_s$ (homogeneous space) does not change ranks.
\end{proof}

\begin{remark}[Status: maximum-principle step under revision per expert review]\label{rem:twopoint-status}
Expert review confirms $m'<0$ on every compact rank-one symmetric space (the Jacobi-field
product form makes $m$ a positive combination of cotangents) and endorses the subordination
frame, but flags three points in the maximum-principle argument that are asserted rather than
proved, all fixable and being repaired: (i) $v=0$ at the antipodal focal set needs proof
(available: each cut stratum has codimension $\ge1$, and smoothness of $p_t$ against the
corner of $\dM$ forces $P'(\diam M)=0$); (ii) the drift $m\sim(d-1)/r$ is singular at the
pole and unbounded at the focal radius, so the strong maximum principle needs Bessel-type
handling or an $\varepsilon$-interior argument; (iii) ``$v<0$ for small $t$'' needs the
uniform small-time asymptotic. The conclusion $\rho_{G,s}=1$ is expected to survive; the
shored-up argument is to be supplied.
\end{remark}

\begin{remark}[What this settles]
Theorem~\ref{thm:twopoint} removes the dimension ceiling from the exact case: the fractional
Green-rank coefficient is $1$ on all rank-one symmetric spaces in \emph{every} dimension, for
every $s>0$---strictly stronger than the $s=1$, $d\le3$ corollary of Paper~I, and proved by a
different (subordination) route that never needs Hilbert--Schmidt stability. Combined with
Theorem~\ref{thm:fraclimit} ($s>d/4$) and Theorem~\ref{thm:fracnear} (openness), the reduction
of Angular Preservation to Green-rank positivity now holds in all dimensions with the
fractional filter, and its exact anchor---value $1$ on the symmetric spaces---is dimension-free.
The residual open problem is the same shape as before, one level up: fractional-Green-rank
positivity $\rho_{G,s}(M)>0$ for general $M$, with the near-radial class of
Theorem~\ref{thm:fracnear} as the proven interior.
\end{remark}



\section{Standing facts, the frontier, and open problems}\label{sec:frontier}

\noindent\textbf{Standing facts.} The proofs draw on the following standard tools, cited at
point of use: the level-set/coarea lemmas and \L ojasiewicz sublevel estimates
(Evans--Gariepy; van den Dries, o-minimality; Comte--Lion--Rolin); cut-locus and heat-kernel
regularity (Sakai; Minakshisundaram--Pleijel; the Green-function parametrix, Aubin); norm
resolvent convergence of uniformly comparable forms (Kato); the uniform pointwise Weyl law
(H\"ormander); the spectral-gap lower bound under a Ricci lower bound and diameter
(Li--Yau; Zhong--Yang); Daniels' inequality and U-statistic consistency for grade
correlations (Daniels~1950; Hoeffding~1948); and the Duhamel heat-kernel perturbation
estimate. None is load-bearing beyond its stated role.

\medskip
\noindent\textbf{What is now reduced to what.} For $d\le3$, Angular Preservation (rank form)
is implied by: (i) the rank transfer, which holds under calibration $E_n\to0$
(resp.\ $E_n\sqrt{A_m}\to0$ at a growing cutoff) and (AC$'$)---the latter proved for
real-analytic metrics (Part~\ref{sec:ac}); and (ii) Green-rank positivity, reduced to a
single scale-invariant non-radiality index (Part~\ref{sec:positivity}), itself the
heat-kernel shell-variance, which vanishes on two-point homogeneous spaces and is controlled
by an explicit spectral-gap tail and one curvature-anisotropy amplitude
(Part~\ref{sec:l2}). Hypotheses (H3) and atomlessness are deleted outright
(Part~\ref{sec:reductions}); the critical dimension is crossed by the fractional filter
(Part~\ref{sec:fractional}).

\medskip
\noindent\textbf{A numerical counterexample hunt.} Green-rank positivity may be false on
extreme geometries; the master inequality gives any counterexample the exact target
$\rho_G\le0$. A registered numerical hunt (public repository,
\texttt{experiments/green-rank}) computes $\rho_G$ across four geometry families---deformed
flat tori, thin-neck surfaces of revolution, Berger spheres, and topological handles---each
with a control validating the discrete pipeline. None yields a counterexample: $\rho_G$
never drops below $\approx0.79$, and the two candidates the theory flags (thin necks and
Berger spheres) both clear. The homogeneous families rise to $\rho_G=1$; inhomogeneity lowers
$\rho_G$ but plateaus, consistent with the shell-variance staying bounded. This is
accumulating evidence, not proof.

\medskip
\noindent\textbf{Open, in order of expected difficulty.} (a) Prove (AC$'$) at fixed cutoff
for generic \emph{smooth} (non-analytic) metrics---a positive-measure level set is the only
obstruction. (b) Make the curvature-anisotropy amplitude of Part~\ref{sec:l2} explicit in
$(K,D,i_0,\sup\lVert\nabla\Ric\rVert)$, promoting the leading-order estimate to a bound
uniform in the head time-interval; this converts the named-$\mathcal A$ class into a closed
curvature/diameter theorem. (c) Green-rank positivity beyond the perturbative and near-radial
regimes---whether the plateau observed numerically reflects a universal lower bound
$\rho_G\ge c>0$ for surfaces, a statement strictly stronger than positivity. (d) The general
$d\ge4$ positivity for the fractional filter---the same question one level up, with the exact
anchor already in hand.

\end{document}
